\section{From Vatican I to the present articulation}

The Reformation dispute intensified questions about Scripture, Tradition,
ecclesial authority, and conscience; confessional polemics sometimes reduced
complex positions on every side. Trent issued dogmatic and reforming decrees
without a single later-style preface assigning a note to every sentence.
Vatican I's \emph{Dei Filius} clarified Revelation, faith, and reason, while
\emph{Pastor aeternus} defined papal primacy and a carefully conditioned papal
infallibility. The council's interruption left fuller episcopal and ecclesial
teaching to later articulation.

Vatican II placed authority within the whole people of God, episcopal
collegiality, Tradition, and the service of the word, especially in
\emph{Lumen gentium} 12 and 25 and \emph{Dei Verbum} 8--10. It did not make
``pastoral'' a synonym for ``optional.'' \emph{Mysterium Ecclesiae} (1973)
defended the permanence of dogmatic truth while acknowledging the historical
conditions of formulations. The 1983 Latin Code juridically articulated
obligations; the 1989 profession, the 1998 commentary, and \emph{Ad tuendam
fidem} aligned the doctrinal and canonical threefold schema. \emph{Donum
veritatis} supplied a moral and ecclesial account of theological vocation and
difficulty.

